\section{TEMA}

O tema deste trabalho é a \textbf{aplicação de Inteligência Artificial Generativa na análise automatizada de falhas em Infraestrutura como Código (IaC)}, especificamente utilizando modelos de linguagem executados localmente para interpretar logs de erro do Terraform e sugerir correções estruturadas.

A escolha deste tema justifica-se pela convergência de duas tendências relevantes na engenharia de software contemporânea: a consolidação da IaC como prática fundamental de DevOps e o amadurecimento dos modelos de linguagem especializados em código. Embora ambas as áreas tenham evoluído significativamente nos últimos anos, a intersecção entre elas --- a aplicação prática de LLMs na resolução de problemas de infraestrutura --- ainda carece de investigação empírica rigorosa, especialmente no contexto de execução local que preserva a privacidade dos dados.

O escopo delimita-se à análise de falhas em código Terraform, utilizando três modelos de linguagem de empresas distintas (Meta, DeepSeek e Alibaba), executados por meio da plataforma Ollama em ambiente local, avaliados em cenários controlados de falhas previamente catalogadas.
